% Imporditud: tools/import_eksperimendid.py
% Allikas: Eesti füüsikaolümpiaadi eksperimendi ülesannete
%   kogu (Rael Kalda, 2023), ülesanne Ü134, lahendus L134.
\setAuthor{EFO žürii}
\setRound{piirkonnavoor}
\setYear{2008}
\setNumber{G E2}
\setDifficulty{4}
\setTopic{dünaamika}
\setVahendid{süstal mahuga 10 ml, joonlaud, baromeeter}
\setPunktid{14}
\setAllikas{Eksperimendi kogu Ü134, lahendus lk 102}
\setLahendusseis{toores}

\prob{Süstal}
Määrata süstla kolvile mõjuv hõõrdejõud. Hinnata mõõteviga.

\hint

\solu
Kolvi pindala saab avaldada kui

S = Vs , l

kus Vs on süstla ruumala kolvi mingi asetuse korral ja l on vastav kolvi käik. Hõõrdejõu saame leida, kui suleme süstla otsa sõrmega jättes sisse mingi koguse gaasi $V_0$ rõhul $p_0$, seejärel surudes gaasi kokku ja lastes gaasil kolb tagasi suruda. Kolb hakkab tekkinud rõhkude vahe mõjul liikuma tagasi kuni hõõrdejõud tasakaalustatakse gaasi rõhuvahest põhjustatud jõu poolt ja seega saame hõõrdejõu avaldada F = S$\Delta$p, kus $\Delta$p on rõhkude erinevus süstla sees ja süstlast väljas.

Eeldades, et temperatuur on konstantne, saame rakendada Boyle-Mariotte'i seadust:

p0V0 = p1V1 $\Rightarrow$ $p_1$ = $p_0$ $V_0$ $\Rightarrow$ $\Delta$p = $p_0$ $V_1$ $-$ $V_0$ = $p_0$$\Delta$V . $V_1$ $V_1$ $V_1$

Ja seega saame hõõrdejõu suuruseks:

F = S$\Delta$p = V $p_0$$\Delta$V . l $V_1$

Mõõtmised võiks teha nii surudes gaas kokku kui ka hõrendades gaasi ja tulemuseks võtta saadud hõõrdejõu keskmise.

Märkus: küsimus, kas vaadeldav protsess süstlas oleva gaasiga on adiabaatiline või isotermiline ei ole triviaalne; vale eelduse kasutamine tekitab 1,4-kordse vea. Õhu difusioonikonstandi abil võib hinnata soojusjuhtivuse tõttu toimuva temperatuuri ühtlustumise karakteerset aega, see tuleb suurusjärgus 1 s. Niisiis tuleks hoolt kanda, et kolvi liigutamine ei toimu väga kiiresti. Antud aspekti analüüsimata jätmise eest punkte alandada pole vaja.

Hindamine: Katse püstituse kirjeldamine --- 2 p. Hõõrdejõu avaldise tuletamine --- 4 p. Mõõtmiste teostamine ja jõu väärtuse arvutamine --- 2 p. Katse korraldamine kahel erineval viisil (gaasi surudes ja hõrendades) ning keskmise võtmine --- 2 p. Tulemus on tõepärane --- 1 p. Mõõtmisi on korratud mitu korda --- 1 p. Mõõtevea hindamine --- 2 p.
\probend
